\lecture{Искусство Киевской Руси}{lecture-02}

{
  \title{История Русского Искусства}
  \subtitle{Лекция 2: Искусство Киевской Руси}
  \date{9 сентября 2026}

  \begin{frame}[plain,noframenumbering]
    \titlepage
  \end{frame}
}

\sectionimage{lectures/02-architecture/images/section-logos/sobor.png}
\section{Архитектура Киевской Руси}

\begin{frame}{Киевская Русь}
  \tallfigure[]
    {lectures/02-architecture/images/maps/map1.png}
    {Местоположение Киевской Руси к середине XI века}

  С принятием христианства в 988~г. Киевское государство было вовлечено
  в культурный поток византийско-славянского мира и восточнохристианскую
  культуру.
\end{frame}

\begin{frame}{Киевская Русь}
    В искусстве этой эпохи преобладало:
    \begin{itemize}
      \item сосредоточение внимания на <<неземном>> / <<невещественном>>;
      \item следование канону (в основном византийскому);
      \item имперсональность. Многие художники, строители, авторы неизвестны.
            Считались исполнителями; роль творца отводилась Церкви.
    \end{itemize}
  \end{frame}

\subsection{Софийский собор в Киеве}

\begin{frame}{Софийский собор}
  \begin{itemize}
    \item Построен сыном Владимира Ярославом Мудрым.
          Летопись относит закладку к 1037~г.
    \item Митрополичий храм, главный собор государства.
    \item Пятинефный крестово-купольный храм с внутренними и внешними
          галереями и двумя лестничными башнями.
    \item Первоначально 13 полусферических куполов.
    \item Мозаики и фрески --- первая половина XI~в. 
    \item Нынешний внешний облик --- украинское барокко конца XVII--XVIII~в.
  \end{itemize}
\end{frame}

\begin{frame}{План}
  \widefigure[П.~А. Раппопорт, \emph{Русская архитектура X--XIII~вв.};
    \href{https://commons.wikimedia.org/wiki/File:Plan_of_Saint_Sofia_Cathedral_in_Kyiv,_Ukraine.jpg}{Wikimedia Commons}.
    CC~BY-SA~4.0]
    {lectures/02-architecture/images/kyiv-sophia/plan.jpg}
    {План Софии Киевской XI~в.: пять нефов, галереи,
     две западные лестничные башни}
\end{frame}

\begin{frame}{Первоначальный облик}
  \widefigure[фото макета: Xristoupolitis;
    \href{https://commons.wikimedia.org/wiki/File:Sofia_Kiev_11th.JPG}{Wikimedia Commons}.
    CC~BY-SA~4.0]
    {lectures/02-architecture/images/kyiv-sophia/reconstruction-11c.jpg}
    {Макет Софии XI~в.: тринадцать полусферических глав,
     открытая плинфяная кладка, двухъярусные галереи}
\end{frame}

\begin{frame}{Восточный фасад}
  \widefigure[фото: Viacheslav Galievskyi;
    \href{https://commons.wikimedia.org/wiki/File:Kyiv-Saint_Sophia_Cathedral-east_view.jpg}{Wikimedia Commons}.
    CC~BY-SA~4.0]
    {lectures/02-architecture/images/kyiv-sophia/exterior-east.jpg}
    {Ряд апсид: под барочной штукатуркой читается
     пятинефный объём XI~века}
\end{frame}

\begin{frame}{Южный фасад}
  \widefigure[фото: Viacheslav Galievskyi;
    \href{https://commons.wikimedia.org/wiki/File:Kyiv-Saint_Sophia_Cathedral-south_view.jpg}{Wikimedia Commons}.
    CC~BY-SA~4.0]
    {lectures/02-architecture/images/kyiv-sophia/exterior-south.jpg}
    {Украинское барокко: надстроенные галереи, контрфорсы,
     грушевидные позолоченные главы}
\end{frame}

\begin{frame}{Интерьер}
  \widefigure[фото: Viacheslav Galievskyi;
    \href{https://commons.wikimedia.org/wiki/File:Kyiv-Saint_Sophia_Cathedral-interior-01.jpg}{Wikimedia Commons}.
    CC~BY-SA~4.0]
    {lectures/02-architecture/images/kyiv-sophia/interior-nave.jpg}
    {Центральный неф: крестово-купольное пространство XI~в.
     Мозаики в куполе и алтаре, фрески на стенах}
\end{frame}

\subsection{Софийские соборы Полоцка и Новгорода}

\begin{frame}{Три Софии}
  Тип киевской Софии переносится в другие центры Руси.
  \begin{itemize}
    \item Полоцк --- Софийский собор Всеслава Брячиславича,
          середина XI~в.
    \item Новгород --- Софийский собор князя Владимира Ярославича,
          1045--1050~гг.
    \item Три храма одного посвящения --- знак политической
          и духовной самостоятельности этих городов.
  \end{itemize}
\end{frame}

\begin{frame}{Полоцкая София}
  \begin{itemize}
    \item Пятинефный крестово-купольный храм на Верхнем замке
          над Западной Двиной.
    \item Первоначально семь шлемовидных глав.
    \item В 1710~г. средневековый собор взорван;
          в середине XVIII~в. И.~К. Глаубиц возводит
          двухбашенную базилику виленского барокко.
    \item От XI~в. остаются фундамент, столбы и восточная апсида.
    \item Мозаики утрачены; фрески XI~в. сохранились лишь фрагменты
          в центральной апсиде.
  \end{itemize}
\end{frame}

\begin{frame}{План}
  \widefigure[И.~Хозеров, 1927;
    \href{https://commons.wikimedia.org/wiki/File:Po\%C5\%82acak,_Vierchni_Zamak,_Sabornaja._\%D0\%9F\%D0\%BE\%D0\%BB\%D0\%B0\%D1\%86\%D0\%B0\%D0\%BA,_\%D0\%92\%D0\%B5\%D1\%80\%D1\%85\%D0\%BD\%D1\%96_\%D0\%97\%D0\%B0\%D0\%BC\%D0\%B0\%D0\%BA,_\%D0\%A1\%D0\%B0\%D0\%B1\%D0\%BE\%D1\%80\%D0\%BD\%D0\%B0\%D1\%8F_(I._Choziera\%C5\%AD,_1927).jpg}{Wikimedia Commons}.
    Общественное достояние]
    {lectures/02-architecture/images/polotsk-sophia/plan.jpg}
    {План первоначальной полоцкой Софии:
     квадрат ок.\ 26$\times$26~м, три гранёные апсиды, без галерей}
\end{frame}

\begin{frame}{Первоначальный облик}
  \tallfigure[фото макета: Da voli;
    \href{https://commons.wikimedia.org/wiki/File:St._Sophia_Cathedral_in_Polotsk_-_Miniature.JPG}{Wikimedia Commons}.
    CC~BY-SA~3.0]
    {lectures/02-architecture/images/polotsk-sophia/reconstruction-model.jpg}
    {Макет Софии XI~в.: крестово-купольный храм
     с семью шлемовидными главами}
\end{frame}

\begin{frame}{Барочный собор}
  \widefigure[фото: Виктор Мелешко;
    \href{https://commons.wikimedia.org/wiki/File:\%D0\%A1\%D0\%BE\%D1\%84\%D0\%B8\%D0\%B9\%D1\%81\%D0\%BA\%D0\%B8\%D0\%B9_\%D1\%81\%D0\%BE\%D0\%B1\%D0\%BE\%D1\%80_\%D0\%B2_\%D0\%9F\%D0\%BE\%D0\%BB\%D0\%BE\%D1\%86\%D0\%BA\%D0\%B5_2.jpg}{Wikimedia Commons}.
    CC~BY-SA~4.0]
    {lectures/02-architecture/images/polotsk-sophia/exterior-baroque.jpg}
    {Нынешний собор Глаубица --- не облик XI~века,
     а барочная базилика на остатках храма Всеслава}
\end{frame}

\begin{frame}{Кладка XI века}
  \widefigure[фото: Чаховіч Уладзіслаў;
    \href{https://commons.wikimedia.org/wiki/File:\%D0\%9F\%D0\%BE\%D0\%BB\%D0\%B0\%D1\%86\%D0\%BA\%D0\%B0\%D1\%8F_\%D0\%A1\%D0\%B0\%D1\%84\%D1\%96\%D1\%8F._\%D0\%A0\%D0\%B0\%D1\%81\%D0\%BA\%D1\%80\%D1\%8B\%D1\%86\%D1\%86\%D1\%96_(01).jpg}{Wikimedia Commons}.
    CC~BY-SA~3.0]
    {lectures/02-architecture/images/polotsk-sophia/masonry-11c.jpg}
    {Раскрытая плинфяная кладка <<утопленным рядом>>
     в нижних частях стен}
\end{frame}

\subsection{Золотые ворота}

\begin{frame}{Золотые ворота}
  \begin{itemize}
    \item Парадный въезд в <<город Ярослава>>, около 1037~г.
    \item Названы по образцу Золотых ворот Константинополя.
    \item Над проездом надвратная церковь Благовещения.
    \item Кладка opus mixtum: плинфа и камень.
    \item К XVII~в. остались руины; в 1982~г. над подлинными стенами
          возведён павильон-реконструкция.
  \end{itemize}
\end{frame}

\begin{frame}{Руины, 1651}
  \widefigure[Абрахам ван Вестерфельд, 1651;
    \href{https://commons.wikimedia.org/wiki/File:The_Golgen_Gate_from_the_west_by_Westerveld.jpg}{Wikimedia Commons}.
    Общественное достояние]
    {lectures/02-architecture/images/golden-gates/westerveld-1651.jpg}
    {Вид с запада: остатки проезда и надвратной церкви;
     на дальнем плане София}
\end{frame}

\begin{frame}{Раскопки}
  \widefigure[литография В.~Тимма, 1858;
    \href{https://commons.wikimedia.org/wiki/File:Zoloti_virita,_1858.jpg}{Wikimedia Commons}.
    Общественное достояние]
    {lectures/02-architecture/images/golden-gates/timm-1858.jpg}
    {Две параллельные стены проезда XI~в.
     после раскопок 1832~г.}
\end{frame}

\begin{frame}{Реконструкция 1982}
  \widefigure[фото: George Chernilevsky, 2018;
    \href{https://commons.wikimedia.org/wiki/File:Golden_Gate_Kiev_2018_G1.jpg}{Wikimedia Commons}.
    Общественное достояние]
    {lectures/02-architecture/images/golden-gates/reconstruction-1982.jpg}
    {Южный фасад павильона: гипотетический облик ворот
     с церковью Благовещения над проездом}
\end{frame}

\begin{frame}{Подлинная кладка}
  \widefigure[фото: Maksym Kozlenko, 2017;
    \href{https://commons.wikimedia.org/wiki/File:2017-05-18_Interior_of_Golden_Gate_(Kiev).jpg}{Wikimedia Commons}.
    CC~BY-SA~4.0]
    {lectures/02-architecture/images/golden-gates/original-masonry.jpg}
    {Внутри павильона --- стены проезда XI~в.
     Это единственные сохранившиеся части ворот Ярослава}
\end{frame}

\subsection{Успенский собор Киево-Печерского монастыря}

\begin{frame}{Успенский собор}
  \begin{itemize}
    \item 1073--1078~гг. Великая церковь Печерского монастыря.
    \item Выразитель усиления христианского аскетизма.
    \item Обширный и высокий трёхнефный храм,
          увенчанный единственным куполом.
    \item Тип, ставший каноном монастырского зодчества Руси,
          в отличие от многоглавой княжеской Софии.
    \item По Патерику работали греческие мастера.
    \item В XVIII~в. барокизирован; взорван в 1941~г.;
          воссоздан в 1990-е--2000 как барочный, не как храм XI~в.
  \end{itemize}
\end{frame}

\begin{frame}{План}
  \tallfigure[\emph{Киевская старина}, 1886;
    \href{https://commons.wikimedia.org/wiki/File:Uspen\%C5\%9Bka_katedra_Kyjevo-Pe\%C4\%8Der\%C5\%9Bkoho_monasty\%C5\%95a._Pozemna_proekcija.png}{Wikimedia Commons}.
    Общественное достояние]
    {lectures/02-architecture/images/pechersk-uspensky/plan-1886.png}
    {Тёмная штриховка --- трёхнефный храм XI~века;
     светлая --- барочные пристройки XVIII~в.}
\end{frame}

\begin{frame}{Первоначальный облик}
  \tallfigure[фото макета: Vi Ko; музей Лавры;
    \href{https://commons.wikimedia.org/wiki/File:The_elements_of_the_Ruins_of_the_Cathedral._03.Model_of_XI_st.jpg}{Wikimedia Commons}.
    CC~BY-SA~4.0]
    {lectures/02-architecture/images/pechersk-uspensky/model-11c.jpg}
    {Макет XI~в.: шестистолпный крестово-купольный храм
     с одним куполом}
\end{frame}

\begin{frame}{До разрушения}
  \widefigure[фото: Гудшон и Губчевский, ок.\ 1912;
    \href{https://commons.wikimedia.org/wiki/File:\%D0\%A3\%D1\%81\%D0\%BF\%D0\%B5\%D0\%BD\%D1\%81\%D1\%8C\%D0\%BA\%D0\%B8\%D0\%B9_\%D1\%81\%D0\%BE\%D0\%B1\%D0\%BE\%D1\%80_\%D0\%BF\%D1\%80\%D0\%B8\%D0\%B1\%D0\%BB\%D0\%B8\%D0\%B7\%D0\%BD\%D0\%BE_1912.jpg}{Wikimedia Commons}.
    Общественное достояние]
    {lectures/02-architecture/images/pechersk-uspensky/historic-1912.jpg}
    {Барочный семиглавый собор. Не аскетичный
     однокупольный храм эпохи Феодосия Печерского}
\end{frame}

\begin{frame}{Воссозданный собор}
  \widefigure[фото: Сергій Венцеславський, 2021;
    \href{https://commons.wikimedia.org/wiki/File:\%D0\%A1\%D0\%BE\%D0\%B1\%D0\%BE\%D1\%80_\%D0\%A3\%D1\%81\%D0\%BF\%D0\%B5\%D0\%BD\%D1\%81\%D1\%8C\%D0\%BA\%D0\%B8\%D0\%B9_2021-2.jpg}{Wikimedia Commons}.
    CC~BY-SA~4.0]
    {lectures/02-architecture/images/pechersk-uspensky/exterior-current.jpg}
    {Освящён в 2000~г. Воспроизводит украинское барокко XVIII~в.}
\end{frame}

\subsection{Княжеские монастыри середины XI века}

\begin{frame}{Два монастыря более раннего времени}
  \begin{itemize}
    \item Дмитриевский монастырь Изяслава Ярославича
          (крестильное имя Дмитрий), середина XI~в.
          На его месте Святополк в 1108--1113~гг. ставит
          Михайловский Златоверхий собор.
    \item Выдубицкий монастырь Всеволода Ярославича.
          Михайловская церковь заложена в 1070~г.,
          освящена в 1088~г.
  \end{itemize}
\end{frame}

\begin{frame}{Выдубичи}
  \widefigure[фото: Galvm;
    \href{https://commons.wikimedia.org/wiki/File:Kyiv-Saint_Michael_Church_(Vydubychi_Monastery)-general_view.jpg}{Wikimedia Commons}.
    CC~BY-SA~4.0]
    {lectures/02-architecture/images/vydubychi/exterior-general.jpg}
    {Михайловская церковь: западная половина XI~в.
     после обрушения алтаря в Днепр и барочной достройки
     1766--1769~гг.}
\end{frame}

\begin{frame}{Кладка XI века}
  \widefigure[фото: Vi Ko;
    \href{https://commons.wikimedia.org/wiki/File:\%D0\%97\%D0\%B0\%D1\%85\%D1\%96\%D0\%B4\%D0\%BD\%D0\%B8\%D0\%B9_\%D1\%84\%D0\%B0\%D1\%81\%D0\%B0\%D0\%B4_\%D1\%81\%D0\%BE\%D0\%B1\%D0\%BE\%D1\%80\%D1\%83.jpg}{Wikimedia Commons}.
    CC~BY-SA~4.0]
    {lectures/02-architecture/images/vydubychi/west-facade.jpg}
    {Западный фасад: нартекс и западные ячейки XI~в.
     с раскрытой плинфяной кладкой}
\end{frame}

\begin{frame}{Михайловский Златоверхий}
  \widefigure[открытка нач.\ XX~в.;
    \href{https://commons.wikimedia.org/wiki/File:Kyiv-Michael-monastery.jpg}{Wikimedia Commons}.
    Общественное достояние]
    {lectures/02-architecture/images/mikhailovsky/historic-c1900.jpg}
    {Собор Святополка на месте Дмитриевского монастыря Изяслава.
     Снесён в 1930-е; нынешний является реконструкцией 1997--1999.
     Мозаики ок.\ 1112~г. сняты до сноса}
\end{frame}

\subsection{Спасо-Преображенский собор в Чернигове}

\begin{frame}{Спас Черниговский}
  \begin{itemize}
    \item Заложен Мстиславом Тмутараканским, братом Ярослава Мудрого.
    \item Начат в 1030-е; Мстислав погребён в недостроенном храме (1036).
          Завершён при Ярославе.
    \item Один из древнейших сохранившихся каменных храмов Руси.
    \item Трёхнефный, с крещальней и лестничной башней.
    \item Кладка плинфой; в интерьере византийские мраморные колонны.
  \end{itemize}
\end{frame}

\begin{frame}{План}
  \widefigure[J.~H. Blasius, 1844;
    \href{https://commons.wikimedia.org/wiki/File:BLASIUS(1844)_p2.634_Grundri\%C3\%9F_der_Christuskirche_von_Tschernigos.jpg}{Wikimedia Commons}.
    Общественное достояние]
    {lectures/02-architecture/images/chernihiv-spaso/plan-1844.jpg}
    {Три нефа, три апсиды, нартекс и круглые западные башни
     (северная --- лестничная XI~в.)}
\end{frame}

\begin{frame}{Внешний вид}
  \widefigure[фото: Nomad0212, 2017;
    \href{https://commons.wikimedia.org/wiki/File:2017_\%D0\%A1\%D0\%BF\%D0\%B0\%D1\%81\%D0\%BE-\%D0\%9F\%D1\%80\%D0\%B5\%D0\%BE\%D0\%B1\%D1\%80\%D0\%B0\%D0\%B6\%D0\%B5\%D0\%BD\%D1\%81\%D1\%8C\%D0\%BA\%D0\%B8\%D0\%B9_\%D1\%81\%D0\%BE\%D0\%B1\%D0\%BE\%D1\%80_\%D1\%83_\%D0\%A7\%D0\%B5\%D1\%80\%D0\%BD\%D1\%96\%D0\%B3\%D0\%BE\%D0\%B2\%D1\%96.jpg}{Wikimedia Commons}.
    CC~BY-SA~4.0]
    {lectures/02-architecture/images/chernihiv-spaso/exterior.jpg}
    {Плинфа, пятиглавие и лестничная башня
     на черниговском Валу}
\end{frame}

\begin{frame}{Апсиды}
  \widefigure[фото: Nomad0212;
    \href{https://commons.wikimedia.org/wiki/File:\%D0\%90\%D0\%BF\%D1\%81\%D0\%B8\%D0\%B4\%D0\%BD\%D0\%B0_\%D1\%87\%D0\%B0\%D1\%81\%D1\%82\%D0\%B8\%D0\%BD\%D0\%B0_\%D0\%A1\%D0\%BF\%D0\%B0\%D1\%81\%D0\%BE-\%D0\%9F\%D1\%80\%D0\%B5\%D0\%BE\%D0\%B1\%D1\%80\%D0\%B0\%D0\%B6\%D0\%B5\%D0\%BD\%D1\%81\%D1\%8C\%D0\%BA\%D0\%BE\%D0\%B3\%D0\%BE_\%D1\%81\%D0\%BE\%D0\%B1\%D0\%BE\%D1\%80\%D1\%83_\%D0\%B2_\%D0\%BC\%D1\%96\%D1\%81\%D1\%82\%D1\%96_\%D0\%A7\%D0\%B5\%D1\%80\%D0\%BD\%D1\%96\%D0\%B3\%D1\%96\%D0\%B2.jpg}{Wikimedia Commons}.
    CC~BY-SA~4.0]
    {lectures/02-architecture/images/chernihiv-spaso/exterior-apses.jpg}
    {Восточный фасад: три апсиды XI~в. и кладка opus mixtum}
\end{frame}

\begin{frame}{Интерьер}
  \tallfigure[фото: Alex Arendar, 2024;
    \href{https://commons.wikimedia.org/wiki/File:\%D0\%86\%D0\%BD\%D1\%82\%D0\%B5\%D1\%80\%27\%D1\%94\%D1\%80_\%D0\%A1\%D0\%BF\%D0\%B0\%D1\%81\%D0\%BE-\%D0\%9F\%D1\%80\%D0\%B5\%D0\%BE\%D0\%B1\%D1\%80\%D0\%B0\%D0\%B6\%D0\%B5\%D0\%BD\%D1\%81\%D1\%8C\%D0\%BA\%D0\%BE\%D0\%B3\%D0\%BE_\%D1\%81\%D0\%BE\%D0\%B1\%D0\%BE\%D1\%80\%D1\%83,\%D0\%B2\%D0\%B8\%D0\%B4_\%D0\%B7_\%D1\%85\%D0\%BE\%D1\%80\%D1\%96\%D0\%B2.jpg}{Wikimedia Commons}.
    CC~BY-SA~4.0]
    {lectures/02-architecture/images/chernihiv-spaso/interior.jpg}
    {Вид с хор: трёхнефное пространство XI~в.
     Росписи и иконостас поздние}
\end{frame}

\begin{frame}{Мраморная колонна}
  \tallfigure[фото: заповедник <<Чернігів стародавній>>;
    \href{https://commons.wikimedia.org/wiki/File:\%D0\%86\%D0\%BD\%D1\%82\%D0\%B5\%D1\%80\%27\%D1\%94\%D1\%80_\%D0\%A1\%D0\%BF\%D0\%B0\%D1\%81\%D0\%B0_\%D0\%A7\%D0\%B5\%D1\%80\%D0\%BD\%D1\%96\%D0\%B3\%D1\%96\%D0\%B2\%D1\%81\%D1\%8C\%D0\%BA\%D0\%BE\%D0\%B3\%D0\%BE.jpg}{Wikimedia Commons}.
    CC0]
    {lectures/02-architecture/images/chernihiv-spaso/marble-column.jpg}
    {Зондаж византийской мраморной колонны с ионийской капителью ---
     подлинная опора XI~в., позднее заключённая в кирпичную обкладку}
\end{frame}

\subsection{Великий Новгород}

\begin{frame}{Новгород на Волхове}
  \begin{itemize}
    \item Город на обоих берегах Волхова.
    \item Левый берег --- Софийская сторона с детинцем.
    \item Правый --- Торговая сторона.
    \item Река --- ось города; собор стоит в северо-западной
          части детинца, над водой.
  \end{itemize}
\end{frame}

\begin{frame}{Детинец}
  \widefigure[фото: SadPetzl;
    \href{https://commons.wikimedia.org/wiki/File:Velikiy_Novgorod_Detinets_aerial_1.jpg}{Wikimedia Commons}.
    CC~BY-SA~4.0]
    {lectures/02-architecture/images/novgorod-city/detinets-aerial.jpg}
    {Аэрофотосъёмка: овальные стены кремля,
     Софийский собор и Волхов}
\end{frame}

\begin{frame}{Вид с реки}
  \widefigure[фото: Anastasiya Lvova;
    \href{https://commons.wikimedia.org/wiki/File:\%D0\%92\%D0\%B8\%D0\%B4_\%D0\%BD\%D0\%B0_\%D0\%9D\%D0\%BE\%D0\%B2\%D0\%B3\%D0\%BE\%D1\%80\%D0\%BE\%D0\%B4\%D1\%81\%D0\%BA\%D0\%B8\%D0\%B9_\%D0\%BA\%D1\%80\%D0\%B5\%D0\%BC\%D0\%BB\%D1\%8C_\%D0\%B8_\%D0\%A1\%D0\%BE\%D1\%84\%D0\%B8\%D0\%B9\%D1\%81\%D0\%BA\%D0\%B8\%D0\%B9_\%D1\%81\%D0\%BE\%D0\%B1\%D0\%BE\%D1\%80_\%D1\%81_\%D0\%92\%D0\%BE\%D0\%BB\%D1\%85\%D0\%BE\%D0\%B2\%D0\%B0.jpg}{Wikimedia Commons}.
    CC~BY-SA~4.0]
    {lectures/02-architecture/images/novgorod-city/from-volkhov.jpg}
    {Стены детинца и шлемовидные главы Софии с Волхова}
\end{frame}

\subsection{Софийский собор в Новгороде}

\begin{frame}{София Новгородская}
  \begin{itemize}
    \item Построена князем Владимиром Ярославичем в 1045--1050~гг.
    \item В центре новгородского детинца.
    \item Пятинефный храм с галереями и лестничной башней.
    \item Местный известняк, а не киевская плинфа.
    \item Пять шлемовидных глав; центральная позолочена в 1408~г.
    \item Один из лучше всего сохранившихся храмов XI~в. на Руси.
  \end{itemize}
\end{frame}

\begin{frame}{План}
  \widefigure[И.~Толстой, Н.~Кондаков, 1899;
    \href{https://commons.wikimedia.org/wiki/File:\%D0\%9F\%D0\%BB\%D0\%B0\%D0\%BD_\%D0\%A1\%D0\%B2\%D1\%8F\%D1\%82\%D0\%BE\%D0\%B9_\%D0\%A1\%D0\%BE\%D1\%84\%D0\%B8\%D0\%B8_\%D0\%9D\%D0\%BE\%D0\%B2\%D0\%B3\%D0\%BE\%D1\%80\%D0\%BE\%D0\%B4\%D1\%81\%D0\%BA\%D0\%BE\%D0\%B9.jpg}{Wikimedia Commons}.
    Общественное достояние]
    {lectures/02-architecture/images/novgorod-sophia/plan.jpg}
    {Пять нефов, три апсиды, галереи, лестничная башня}
\end{frame}

\begin{frame}{В детинце}
  \widefigure[фото: Dennis G. Jarvis;
    \href{https://commons.wikimedia.org/wiki/File:Russia_3207_-_St._Sophia_Cathedral_(4138551353).jpg}{Wikimedia Commons}.
    CC~BY-SA~2.0]
    {lectures/02-architecture/images/novgorod-sophia/exterior-detinets.jpg}
    {Известняковые стены, свинцовые боковые главы
     и золочёный центральный купол}
\end{frame}

\begin{frame}{Восточный фасад}
  \widefigure[фото: Kamenskiyan;
    \href{https://commons.wikimedia.org/wiki/File:\%D0\%A1\%D0\%BE\%D1\%84\%D0\%B8\%D0\%B9\%D1\%81\%D0\%BA\%D0\%B8\%D0\%B9_\%D0\%A1\%D0\%BE\%D0\%B1\%D0\%BE\%D1\%80_\%D0\%B2_\%D0\%94\%D0\%B5\%D1\%82\%D0\%B8\%D0\%BD\%D1\%86\%D0\%B5_(\%D0\%B2\%D0\%B8\%D0\%B4_\%D1\%81_\%D0\%B2\%D0\%BE\%D1\%81\%D1\%82\%D0\%BE\%D0\%BA\%D0\%B0).jpg}{Wikimedia Commons}.
    CC~BY-SA~4.0]
    {lectures/02-architecture/images/novgorod-sophia/exterior-east.jpg}
    {Трёхапсидный объём и вертикальный силуэт пяти глав}
\end{frame}

\begin{frame}{Главы}
  \widefigure[фото: Buzov Vladimir;
    \href{https://commons.wikimedia.org/wiki/File:\%D0\%9A\%D1\%83\%D0\%BF\%D0\%BE\%D0\%BB\%D0\%B0_\%D0\%A1\%D0\%BE\%D1\%84\%D0\%B8\%D0\%B9\%D1\%81\%D0\%BA\%D0\%BE\%D0\%B3\%D0\%BE_\%D1\%81\%D0\%BE\%D0\%B1\%D0\%BE\%D1\%80\%D0\%B0,_\%D0\%92\%D0\%B5\%D0\%BB\%D0\%B8\%D0\%BA\%D0\%B8\%D0\%B9_\%D0\%9D\%D0\%BE\%D0\%B2\%D0\%B3\%D0\%BE\%D1\%80\%D0\%BE\%D0\%B4.jpg}{Wikimedia Commons}.
    CC~BY-SA~4.0]
    {lectures/02-architecture/images/novgorod-sophia/helmet-domes.jpg}
    {Четыре свинцовых шлема и золочёный центральный купол}
\end{frame}

\begin{frame}{Магдебургские врата}
  \tallfigure[фото: Александр Сигачёв, 2025;
    \href{https://commons.wikimedia.org/wiki/File:Velikiy_Novgorod_-_P\%C5\%82ock_Door_-_2025-08_-_p1.jpg}{Wikimedia Commons}.
    CC~BY-SA~4.0]
    {lectures/02-architecture/images/novgorod-sophia/magdeburg-doors.jpg}
    {Бронзовые врата середины XII~в. на западном портале
     (Плоцкие / Магдебургские). Рельефы библейских сцен}
\end{frame}

\begin{frame}{Интерьер}
  \widefigure[фото: Дар Ветер;
    \href{https://commons.wikimedia.org/wiki/File:Sophia_inside.jpg}{Wikimedia Commons}.
    CC~BY-SA~3.0]
    {lectures/02-architecture/images/novgorod-sophia/interior.jpg}
    {Высокие столпы и узкие компартименты ---
     более вертикальный объём, чем у Софии Киевской}
\end{frame}

\begin{frame}{Константин и Елена}
  \tallfigure[фото: Lumaca;
    \href{https://commons.wikimedia.org/wiki/File:\%D0\%9A\%D0\%BE\%D0\%BD\%D1\%81\%D1\%82\%D0\%B0\%D0\%BD\%D1\%82\%D0\%B8\%D0\%BD_\%D0\%B8_\%D0\%95\%D0\%BB\%D0\%B5\%D0\%BD\%D0\%B0_(\%D0\%A1\%D0\%BE\%D1\%84\%D0\%B8\%D0\%B9\%D1\%81\%D0\%BA\%D0\%B8\%D0\%B9_\%D1\%81\%D0\%BE\%D0\%B1\%D0\%BE\%D1\%80).jpg}{Wikimedia Commons}.
    CC~BY-SA~4.0]
    {lectures/02-architecture/images/novgorod-sophia/fresco-constantine-helena.jpg}
    {Один из древнейших фрагментов росписи собора, XI~в.
     У южного входа; связан с освящением храма}
\end{frame}

\begin{frame}{Пётр и Павел}
  \tallfigure[Новгородский музей;
    \href{https://commons.wikimedia.org/wiki/File:Ancient_icon_of_sts_peter_\%26_paul.jpg}{Wikimedia Commons}.
    Общественное достояние]
    {lectures/02-architecture/images/novgorod-sophia/icon-peter-paul.jpg}
    {Икона апостолов Петра и Павла, ок.\ 1050~г.,
     из Софийского собора. Одна из древнейших икон,
     связанных с самим храмом}
\end{frame}

\sectionimage{sirin.png}
\section{Живопись Киевской Руси}

\begin{frame}{Мозаика и фреска}
  \begin{itemize}
    \item Мозаика --- купол, паруса, алтарь: золотой фон,
          литургическая иерархия.
    \item Фреска --- стены, столпы, башни: повествовательные циклы.
    \item Мастера --- византийские и местные; имена неизвестны.
  \end{itemize}
\end{frame}

\subsection{Мозаики Софии Киевской}

\begin{frame}{Программа купола}
  Сверху вниз:
  \begin{itemize}
    \item Вседержитель в медальоне, в окружении четырёх архангелов;
    \item апостолы в простенках между окнами барабана;
    \item евангелисты на парусах;
    \item Благовещение на восточных столбах
          подкупольного квадрата;
    \item в конхе апсиды --- Богоматерь Оранта.
  \end{itemize}
\end{frame}

\begin{frame}{Вседержитель}
  \widefigure[фото: Rasal Hague;
    \href{https://commons.wikimedia.org/wiki/File:\%D0\%9F\%D0\%B0\%D0\%BD\%D1\%82\%D0\%BE\%D0\%BA\%D1\%80\%D0\%B0\%D1\%82\%D0\%BE\%D1\%80_(20210903).jpg}{Wikimedia Commons}.
    CC~BY-SA~4.0]
    {lectures/02-architecture/images/kyiv-sophia/mosaic-pantocrator.jpg}
    {Христос Пантократор в медальоне центрального купола.
     Один из архангелов --- подлинная мозаика XI~в.;
     остальные дописаны М.~Врубелем}
\end{frame}

\begin{frame}{Барабан и паруса}
  \widefigure[фото: Suicasmo, 2019;
    \href{https://commons.wikimedia.org/wiki/File:Interior_of_Saint_Sophia_Cathedral_20190505.jpg}{Wikimedia Commons}.
    CC~BY-SA~4.0]
    {lectures/02-architecture/images/sophia-mosaics/dome-drum-apostles.jpg}
    {Апостолы в простенках между окнами барабана
     и евангелисты на парусах.
     Кроме Павла фигуры апостолов восполнены маслом XIX~в.}
\end{frame}

\begin{frame}{Апостол Павел}
  \tallfigure[Национальный заповедник <<София Киевская>>;
    \href{https://commons.wikimedia.org/wiki/File:The_Apostle_Paul_(detail)_-_Google_Art_Project.jpg}{Wikimedia Commons}.
    Общественное достояние]
    {lectures/02-architecture/images/sophia-mosaics/apostle-paul.jpg}
    {Мозаика апостола Павла в простенке барабана ---
     единственный сохранившийся подлинный апостол XI~в.}
\end{frame}

\begin{frame}{Евангелист Марк}
  \tallfigure[Национальный заповедник <<София Киевская>>;
    \href{https://commons.wikimedia.org/wiki/File:St._Mark_the_Evangelist_-_Google_Art_Project.jpg}{Wikimedia Commons}.
    Общественное достояние]
    {lectures/02-architecture/images/sophia-mosaics/evangelist-mark.jpg}
    {Мозаика на парусе --- единственное полностью
     сохранившееся изображение евангелиста.
     Иоанн --- фрагмент; Матфей и Лука --- масло XIX~в.}
\end{frame}

\begin{frame}{Благовещение. Гавриил}
  \tallfigure[Национальный заповедник <<София Киевская>>;
    \href{https://commons.wikimedia.org/wiki/File:The_Annunciation._The_Archangel_Gabriel_-_Google_Art_Project.jpg}{Wikimedia Commons}.
    Общественное достояние]
    {lectures/02-architecture/images/sophia-mosaics/annunciation-gabriel.jpg}
    {Архангел Гавриил на северном восточном столбе
     подкупольного квадрата}
\end{frame}

\begin{frame}{Благовещение. Богоматерь}
  \tallfigure[Национальный заповедник <<София Киевская>>;
    \href{https://commons.wikimedia.org/wiki/File:The_Annunciation._The_Virgin_Mary_-_Google_Art_Project.jpg}{Wikimedia Commons}.
    Общественное достояние]
    {lectures/02-architecture/images/sophia-mosaics/annunciation-virgin.jpg}
    {Богоматерь с пряжей на южном восточном столбе.
     Два образа читаются навстречу друг другу через пространство нефа}
\end{frame}

\begin{frame}{Оранта}
  \tallfigure[Национальный заповедник <<София Киевская>>;
    \href{https://commons.wikimedia.org/wiki/File:Oranta-Kyiv.jpg}{Wikimedia Commons}.
    Общественное достояние]
    {lectures/02-architecture/images/kyiv-sophia/mosaic-oranta.jpg}
    {Богоматерь Оранта (<<Нерушимая Стена>>) в конхе главной апсиды.
     Высота ок.\ 6~м. Первая половина XI~в.}
\end{frame}

\subsection{Фрески софийских храмов}

\begin{frame}{Фресковые циклы}
  Стены Софии Киевской покрыты фресками XI~в.:
  \begin{itemize}
    \item сцены земной жизни Христа;
    \item житие Богоматери (придел Иоакима и Анны);
    \item цикл архангела Михаила (придел и лестничная башня);
    \item ктиторский портрет семьи Ярослава Мудрого.
  \end{itemize}
  В новгородской Софии от XI~в. сохранились лишь фрагменты
  (Константин и Елена --- выше).
\end{frame}

\begin{frame}{Сошествие во ад}
  \widefigure[Национальный заповедник <<София Киевская>>;
    \href{https://commons.wikimedia.org/wiki/File:The_Descent_of_Christ_into_Hell_-_Google_Art_Project.jpg}{Wikimedia Commons}.
    Общественное достояние]
    {lectures/02-architecture/images/sophia-frescoes/anastasis.jpg}
    {Анастасис --- сцена из христологического цикла, XI~в.}
\end{frame}

\begin{frame}{Уверение Фомы}
  \widefigure[Национальный заповедник <<София Киевская>>;
    \href{https://commons.wikimedia.org/wiki/File:The_Doubting_of_Thomas_-_Google_Art_Project.jpg}{Wikimedia Commons}.
    Общественное достояние]
    {lectures/02-architecture/images/sophia-frescoes/doubting-thomas.jpg}
    {Сцена из цикла земной жизни Христа}
\end{frame}

\begin{frame}{Рождество Богородицы}
  \tallfigure[Национальный заповедник <<София Киевская>>;
    \href{https://commons.wikimedia.org/wiki/File:The_Nativity_of_the_Virgin_-_Google_Art_Project.jpg}{Wikimedia Commons}.
    Общественное достояние]
    {lectures/02-architecture/images/sophia-frescoes/nativity-virgin.jpg}
    {Фреска богородичного цикла в приделе Иоакима и Анны}
\end{frame}

\begin{frame}{Целование Марии и Елисаветы}
  \tallfigure[Национальный заповедник <<София Киевская>>;
    \href{https://commons.wikimedia.org/wiki/File:The_Visitation_-_Google_Art_Project.jpg}{Wikimedia Commons}.
    Общественное достояние]
    {lectures/02-architecture/images/sophia-frescoes/visitation.jpg}
    {Ещё одна сцена жития Богоматери, XI~в.}
\end{frame}

\begin{frame}{Архангел Михаил}
  \widefigure[Национальный заповедник <<София Киевская>>;
    \href{https://commons.wikimedia.org/wiki/File:The_Archangel_Michael_-_Google_Art_Project_(LwE83yRr6_HWQQ).jpg}{Wikimedia Commons}.
    Общественное достояние]
    {lectures/02-architecture/images/sophia-frescoes/archangel-michael.jpg}
    {Поясной образ в конхе Михайловского придела Софии Киевской}
\end{frame}

\begin{frame}{Явление архангела Валааму}
  \widefigure[Национальный заповедник <<София Киевская>>;
    \href{https://commons.wikimedia.org/wiki/File:The_Appearance_of_the_Archangel_to_Balaam_-_Google_Art_Project.jpg}{Wikimedia Commons}.
    Общественное достояние]
    {lectures/02-architecture/images/sophia-frescoes/archangel-balaam.jpg}
    {Сцена михайловского цикла в лестничной башне}
\end{frame}

\begin{frame}{Семья Ярослава}
  \widefigure[Национальный заповедник <<София Киевская>>;
    \href{https://commons.wikimedia.org/wiki/File:Princely_group_portrait._South_wall_of_the_nave._-_Google_Art_Project.jpg}{Wikimedia Commons}.
    Общественное достояние]
    {lectures/02-architecture/images/kyiv-sophia/fresco-yaroslav-family.jpg}
    {Сохранившийся фрагмент ктиторской фрески на южной стене:
     члены княжеской семьи (традиционно --- дети Ярослава)}
\end{frame}

\begin{frame}{Полная композиция}
  \widefigure[Абрахам ван Вестерфельд, 1651;
    \href{https://commons.wikimedia.org/wiki/File:Yaroslav_the_Wise_and_his_family_(fresco)_by_Westerveld.jpg}{Wikimedia Commons}.
    Общественное достояние]
    {lectures/02-architecture/images/sophia-frescoes/yaroslav-family-westerveld.jpg}
    {Рисунок полной ктиторской композиции:
     Ярослав подносит Христу модель собора.
     Центральная часть на западной стене разобрана в XVII~в.}
\end{frame}

\subsection{Мозаики Михайловского Златоверхого}

\begin{frame}{Михайловский Златоверхий, ок.\ 1112}
  \begin{itemize}
    \item Собор Святополка Изяславича на месте Дмитриевского
          монастыря середины XI~в.
    \item Мозаики ок.\ 1112~г. --- уже местная школа
          при византийском каноне.
    \item В 1934~г. сняты до сноса храма.
    \item Евхаристия --- в заповеднике <<София Киевская>>;
          Дмитрий Солунский --- в Третьяковской галерее.
  \end{itemize}
\end{frame}

\begin{frame}{Евхаристия}
  \widefigure[фото: Rasal Hague; заповедник <<София Киевская>>;
    \href{https://commons.wikimedia.org/wiki/File:\%D0\%9C\%D0\%BE\%D0\%B7\%D0\%B0\%D1\%97\%D0\%BA\%D0\%B8_\%D0\%97\%D0\%BE\%D0\%BB\%D0\%BE\%D1\%82\%D0\%BE\%D0\%B2\%D0\%B5\%D1\%80\%D1\%85\%D0\%BE\%D0\%B3\%D0\%BE_\%D1\%81\%D0\%BE\%D0\%B1\%D0\%BE\%D1\%80\%D1\%83.jpg}{Wikimedia Commons}.
    CC~BY-SA~4.0]
    {lectures/02-architecture/images/mikhailovsky/mosaic-eucharist.jpg}
    {Причащение апостолов из алтаря Михайловского собора.
     Не путать с Евхаристией Софии Киевской}
\end{frame}

\begin{frame}{Дмитрий Солунский}
  \tallfigure[Третьяковская галерея;
    \href{https://commons.wikimedia.org/wiki/File:\%D0\%A1\%D0\%B2._\%D0\%94\%D0\%BC\%D0\%B8\%D1\%82\%D1\%80\%D0\%B8\%D0\%B9_\%D0\%A1\%D0\%BE\%D0\%BB\%D1\%83\%D0\%BD\%D1\%81\%D0\%BA\%D0\%B8\%D0\%B9_-_Google_Art_Project.jpg}{Wikimedia Commons}.
    Общественное достояние]
    {lectures/02-architecture/images/mikhailovsky/mosaic-demetrius.jpg}
    {Мозаика со столба Михайловского собора ---
     патрональный образ Изяслава-Дмитрия}
\end{frame}

\sectionimage{lectures/02-architecture/images/section-logos/book.png}
\section{Книжная миниатюра Киевской Руси}

\begin{frame}{Книжная миниатюра}
  \begin{itemize}
    \item Рукопись --- княжеский заказ, как и храм.
    \item Пергамен, устав, киноварь, золото.
    \item Два главных кодекса XI~в.: Остромирово евангелие
          (1056--1057) и Изборник Святослава 1073~г.
  \end{itemize}
\end{frame}

\subsection{Остромирово евангелие}

\begin{frame}{Остромирово евангелие}
  \begin{itemize}
    \item 1056--1057~гг. Диакон Григорий.
    \item Для новгородского посадника Остромира,
          свойственника князя Изяслава Ярославича.
    \item Древнейшая датированная книга Руси.
    \item Три миниатюры евангелистов: Иоанн, Лука, Марк
          (Матфей утрачен).
  \end{itemize}
\end{frame}

\begin{frame}{Евангелист Марк}
  \tallfigure[РНБ, F.п.I.5, л.~126;
    \href{https://commons.wikimedia.org/wiki/File:Ostromir_Gospel.jpg}{Wikimedia Commons}.
    Общественное достояние]
    {lectures/02-architecture/images/ostromir-gospel/evangelist-mark.jpg}
    {Марк в квадрифолии; символ --- лев.
     Остромирово евангелие, 1056--1057}
\end{frame}

\begin{frame}{Евангелист Лука}
  \tallfigure[РНБ, F.п.I.5, л.~87~об.;
    \href{https://commons.wikimedia.org/wiki/File:Ostromir_luke.jpg}{Wikimedia Commons}.
    Общественное достояние]
    {lectures/02-architecture/images/ostromir-gospel/evangelist-luke.jpg}
    {Лука с символом тельца.
     Остромирово евангелие, 1056--1057}
\end{frame}

\subsection{Изборник Святослава 1073}

\begin{frame}{Изборник 1073}
  \begin{itemize}
    \item Роскошная киевская рукопись, заказанная
          Святославом Ярославичем.
    \item ГИМ, Син.~1043.
    \item Не путать с Изборником 1076~г.
    \item Миниатюры: княжеская семья; Спас на престоле;
          фронтисписы с отцами церкви; эмальерные заставки.
  \end{itemize}
\end{frame}

\begin{frame}{Семья Святослава}
  \tallfigure[ГИМ, Син.~1043, л.~1~об.;
    \href{https://commons.wikimedia.org/wiki/File:\%D0\%98\%D0\%B7\%D0\%B1\%D0\%BE\%D1\%80\%D0\%BD\%D0\%B8\%D0\%BA_\%D0\%A1\%D0\%B2\%D1\%8F\%D1\%82\%D0\%BE\%D1\%81\%D0\%BB\%D0\%B0\%D0\%B2\%D0\%B0,_XI_\%D0\%B2\%D0\%B5\%D0\%BA._\%D0\%9A\%D0\%BD\%D1\%8F\%D0\%B6\%D0\%B5\%D1\%81\%D0\%BA\%D0\%B0\%D1\%8F_\%D1\%81\%D0\%B5\%D0\%BC\%D1\%8C\%D1\%8F.png}{Wikimedia Commons}.
    Общественное достояние]
    {lectures/02-architecture/images/izbornik-1073/family-portrait.png}
    {Святослав Ярославич с княгиней и сыновьями ---
     княжеский ктиторский портрет в книге}
\end{frame}

\begin{frame}{Разворот}
  \widefigure[фото: Shakko; ГИМ;
    \href{https://commons.wikimedia.org/wiki/File:Svyatoslav\%27s_Miscellanies_03_by_shakko.jpg}{Wikimedia Commons}.
    CC~BY-SA~4.0]
    {lectures/02-architecture/images/izbornik-1073/opening-savior.jpg}
    {Оригинал: семья (л.~1~об.) и Спас на престоле (л.~2).
     Лист со Спасом сильно потёрт}
\end{frame}

\begin{frame}{Фронтиспис}
  \tallfigure[ГИМ, Син.~1043;
    \href{https://commons.wikimedia.org/wiki/File:Izbornyk_1073.jpg}{Wikimedia Commons}.
    Общественное достояние]
    {lectures/02-architecture/images/izbornik-1073/frontispiece-fathers.jpg}
    {<<Храм>> с отцами церкви, павлинами
     и орнаментом эмальерного типа}
\end{frame}

\sectionimage{lectures/02-architecture/images/section-logos/applied_art.png}
\section{Скульптура и прикладное искусство Киевской Руси}

\begin{frame}{Прикладное искусство}
  \begin{itemize}
    \item Украшает княжеский быт так же, как мозаика --- храм.
    \item Каменная резьба на Руси XI~в. редка.
    \item Главные техники металла: перегородчатая эмаль,
          чеканка, скань, ажурное литьё.
  \end{itemize}
\end{frame}

\subsection{Саркофаг Ярослава Мудрого}

\begin{frame}{Саркофаг Ярослава}
  \begin{itemize}
    \item Мраморная рака в северной галерее Софии Киевской.
    \item Традиционно --- гробница Ярослава (ум.\ 1054).
    \item Проконесский мрамор; барельефы: лоза, пальмы,
          птицы, рыбы, кресты.
    \item Вероятно, византийская работа, привезённая в Киев.
  \end{itemize}
\end{frame}

\begin{frame}{Саркофаг}
  \widefigure[фото: Petar Milošević; София Киевская;
    \href{https://commons.wikimedia.org/wiki/File:\%D0\%A1\%D0\%B0\%D1\%80\%D0\%BA\%D0\%BE\%D1\%84\%D0\%B0\%D0\%B3_\%D0\%AF\%D1\%80\%D0\%BE\%D1\%81\%D0\%BB\%D0\%B0\%D0\%B2\%D0\%B0_\%D0\%9C\%D1\%83\%D0\%B4\%D1\%80\%D0\%BE\%D0\%B3\%D0\%BE.jpg}{Wikimedia Commons}.
    CC~BY-SA~3.0]
    {lectures/02-architecture/images/yaroslav-sarcophagus/sarcophagus.jpg}
    {Мраморный саркофаг в Софии Киевской}
\end{frame}

\begin{frame}{Резьба}
  \widefigure[фото: София Киевская;
    \href{https://commons.wikimedia.org/wiki/File:\%D0\%A1\%D0\%B0\%D1\%80\%D0\%BA\%D0\%BE\%D1\%84\%D0\%B0\%D0\%B3_\%D0\%AF\%D1\%80\%D0\%BE\%D1\%81\%D0\%BB\%D0\%B0\%D0\%B2\%D0\%B0.jpg}{Wikimedia Commons}.
    CC~BY-SA~4.0]
    {lectures/02-architecture/images/yaroslav-sarcophagus/sarcophagus-angle.jpg}
    {Двускатная крышка с акротериями и резные стенки}
\end{frame}

\begin{frame}{Орнамент}
  \widefigure[рисунок М.~Берлинского, 1820;
    \href{https://commons.wikimedia.org/wiki/File:Tomb_of_Yaroslav_I_the_Wise_drawing.jpg}{Wikimedia Commons}.
    Общественное достояние]
    {lectures/02-architecture/images/yaroslav-sarcophagus/sarcophagus-drawing.jpg}
    {Программа барельефа: лоза, пальмы, христианские символы}
\end{frame}

\subsection{Украшенная книга}

\begin{frame}{Книга как предмет роскоши}
  Помимо миниатюр --- инициалы, заставки, каллиграфия устава.
  Те же кодексы, что и в предыдущем разделе.
\end{frame}

\begin{frame}{Заставка}
  \tallfigure[РНБ, F.п.I.5, л.~2;
    \href{https://commons.wikimedia.org/wiki/File:Ostromir_Gospel_1.jpg}{Wikimedia Commons}.
    Общественное достояние]
    {lectures/02-architecture/images/ostromir-gospel/zastavka.jpg}
    {Заставка и инициал Остромирова евангелия, 1056--1057}
\end{frame}

\begin{frame}{Инициал и устав}
  \tallfigure[фото: Shakko, 2011;
    \href{https://commons.wikimedia.org/wiki/File:Ostromir_Gospel_4.jpg}{Wikimedia Commons}.
    CC~BY-SA~3.0]
    {lectures/02-architecture/images/ostromir-gospel/folio-initial.jpg}
    {Двустолпный устав и крупный инициал ---
     Остромирово евангелие, оригинал}
\end{frame}

\begin{frame}{Заставка Изборника}
  \tallfigure[ГИМ, Син.~1043;
    \href{https://commons.wikimedia.org/wiki/File:Izbornik_page.jpg}{Wikimedia Commons}.
    Общественное достояние]
    {lectures/02-architecture/images/izbornik-1073/headpiece.jpg}
    {Квадратная заставка эмальерного типа
     и начало текста в два столбца}
\end{frame}

\subsection{Златокузнецы и перегородчатая эмаль}

\begin{frame}{Киевская эмаль}
  \begin{itemize}
    \item Перегородчатая эмаль --- византийская техника,
          освоенная киевскими мастерскими XI--XII~вв.
    \item Колты, рясна, диадемы из княжеских кладов.
    \item Находки у Десятинной церкви (Киев),
          в Чернигове, в Старой Рязани --- до 1237~г.
  \end{itemize}
\end{frame}

\begin{frame}{Клад у Десятинной}
  \widefigure[фото: shakko; ГИМ;
    \href{https://commons.wikimedia.org/wiki/File:Klad_desyatinnaya.jpg}{Wikimedia Commons}.
    CC~BY-SA~3.0]
    {lectures/02-architecture/images/applied-art/kolty-ryasna-tithes.jpg}
    {Золотые колты и рясна с перегородчатой эмалью.
     Клад у Десятинной церкви, Киев, 1-я пол.\ XII~в.}
\end{frame}

\begin{frame}{Рясна}
  \widefigure[фото: Лапоть; ГИМ;
    \href{https://commons.wikimedia.org/wiki/File:Ryasna_Kiev_12c_GIM_Zlatnik.jpg}{Wikimedia Commons}.
    CC0]
    {lectures/02-architecture/images/applied-art/ryasna-tithes.jpg}
    {Золотая рясна с эмалевыми медальонами.
     Тот же киевский клад, XII~в.}
\end{frame}

\begin{frame}{Колты из Чернигова}
  \widefigure[фото: shakko; Эрмитаж;
    \href{https://commons.wikimedia.org/wiki/File:Kolty_(Hermitage)_01.JPG}{Wikimedia Commons}.
    CC~BY-SA~3.0]
    {lectures/02-architecture/images/applied-art/kolty-chernihiv.jpg}
    {Золотые колты с перегородчатой эмалью.
     Клад 1887~г., Александровская площадь, Чернигов; XII~в.}
\end{frame}

\begin{frame}{Диадема}
  \widefigure[фото: shakko; ГРМ;
    \href{https://commons.wikimedia.org/wiki/File:Diademe_(Kiev,12th_c.,_GRM)_by_shakko.jpg}{Wikimedia Commons}.
    CC~BY-SA~3.0]
    {lectures/02-architecture/images/applied-art/diadem-kiev.jpg}
    {Золотая диадема с жемчугом и эмалью.
     Киевский клад 1889~г., XI--XII~вв.}
\end{frame}

\begin{frame}{Браслет}
  \tallfigure[клад Старой Рязани;
    \href{https://commons.wikimedia.org/wiki/File:Bracelet_Razan.jpg}{Wikimedia Commons}.
    CC0]
    {lectures/02-architecture/images/applied-art/bracelet-ryazan.jpg}
    {Серебряный браслет, 2-я пол.\ XII~в.
     Старая Рязань --- домонгольская Русь}
\end{frame}

\begin{frame}{Клад Старой Рязани}
  \tallfigure[фото: shakko; ГИМ;
    \href{https://commons.wikimedia.org/wiki/File:Staraya_Ryazan_treasure_(13th_c.,_GIM)_by_shakko_01.jpg}{Wikimedia Commons}.
    CC~BY-SA~4.0]
    {lectures/02-architecture/images/applied-art/staraya-ryazan-treasure.jpg}
    {Колты, рясна, височные кольца.
     Городище Старая Рязань, 1-я четв.\ XIII~в., до 1237~г.}
\end{frame}

\subsection{Посуда и парадное оружие}

\begin{frame}{Дворцовая утварь}
  \begin{itemize}
    \item Золотая и серебряная посуда княжеских палат
          почти не дошла: переплавка, пожары, 1240~г.
    \item Свободной достоверной фотографии княжеской чаши XI--XII~в.
          нет (чара Владимира Давыдовича --- надпись спорной даты).
    \item Парадное оружие и наборные пояса той же участи:
          известны по кладам, но без устойчивых свободных снимков.
    \item Представление о дворцовой роскоши дают ювелирные клады:
          колты, рясна, диадемы, браслеты.
  \end{itemize}
\end{frame}